\documentclass[11pt,a4paper]{article}
\usepackage[spanish,es-noshorthands,es-tabla]{babel}
\usepackage[T1]{fontenc}
\usepackage[utf8]{inputenc}
\usepackage[defaultfam]{montserrat}
\usepackage{graphicx,tabularx,array,xcolor,colortbl,geometry,fancyhdr,enumitem,needspace}
\geometry{left=2cm,right=2cm,top=2.1cm,bottom=1.8cm}
\definecolor{LBcian}{HTML}{016D84}
\definecolor{LBcereza}{HTML}{A3195B}
\definecolor{LBgris}{HTML}{EFEFEF}
\definecolor{LBtexto}{HTML}{333333}
\setlength{\parindent}{0pt}
\setlength{\headheight}{15pt}
\setlength{\tabcolsep}{6pt}
\renewcommand{\arraystretch}{1.35}
\pagestyle{fancy}\fancyhf{}
\fancyhead[L]{\small\color{LBcian}\bfseries IES LUIS BRAILLE}
\fancyhead[R]{\small Bases de Datos \textperiodcentered{} RA3}
\fancyfoot[L]{\small Departamento de Informática \textperiodcentered{} Coslada}
\fancyfoot[R]{\small Página \thepage}
\renewcommand{\headrulewidth}{0.4pt}
\renewcommand{\footrulewidth}{0pt}
\begin{document}
% Recursos institucionales de la plantilla del repositorio.
\noindent\begin{tabularx}{\textwidth}{@{}Xcc@{}}
  \begin{minipage}[c]{0.55\textwidth}
    {\color{LBcian}\Large\bfseries IES LUIS BRAILLE}\\[2pt]
    {\color{LBtexto}\small Departamento de Informática \textperiodcentered{} Coslada}
  \end{minipage}
  & \includegraphics[width=1.45cm,height=1.3cm,keepaspectratio]{../../../Recursos/Examenes/PlantillaExamenLaTeX/img/logoUE.jpg}
  & \includegraphics[width=3.0cm,height=1.3cm,keepaspectratio]{../../../Recursos/Examenes/PlantillaExamenLaTeX/img/logoCAM.png}
\end{tabularx}
\vspace{2mm}\color{LBcian}\rule{\textwidth}{1.2pt}\color{LBtexto}
\vspace{1mm}
\begin{center}
  {\color{LBcian}\Large\bfseries Examen de consultas SQL \textperiodcentered{} Modelo C}\\[2mm]
  {\color{LBcereza}\small\bfseries BASES DE DATOS \textperiodcentered{} 1.º DAW \textperiodcentered{} RA3}
\end{center}
\noindent\begin{tabularx}{\textwidth}{|>{\bfseries}p{2.9cm}|X|>{\bfseries}p{2.4cm}|X|}\hline
Módulo & Bases de Datos & Curso & 1.º DAW \\ \hline
Nombre & & Apellidos & \\ \hline
Fecha & & Calificación & \cellcolor{LBgris}\hspace{1.2cm}/ 10 \\ \hline
Base de datos & Sakila (MySQL) & Evaluación & RA3 \\ \hline
\end{tabularx}
\medskip\noindent\begin{tabularx}{\textwidth}{|X|}\hline
\cellcolor{LBgris}\color{LBcian}\bfseries INSTRUCCIONES \\ \hline
\begin{minipage}{\dimexpr\linewidth-2\tabcolsep\relax}\vspace{2mm}\small \begin{itemize}[leftmargin=4mm,nosep]\item Escribe una \textbf{consulta SQL independiente} para cada ejercicio. Cada enunciado describe las columnas, las filas y la ordenación que debe mostrar el resultado.
\item El alumnado puede traer en papel la hoja de consulta (cheat sheet) y debe traer también en papel el diagrama de \mbox{Workbench}. Se permiten anotaciones en ambos documentos.
\item Se valorará la \textbf{corrección lógica} y una sintaxis \textbf{razonablemente correcta}. Un error menor de sintaxis no anula toda la puntuación si se entiende claramente la solución propuesta.
\item Se admitirán distintas soluciones que produzcan el resultado solicitado. El examen se limita a consultas sobre una tabla y relaciones internas entre tablas: no se admiten uniones externas, productos cartesianos, subconsultas, vistas ni CTE.
\end{itemize}\vspace{2mm}\end{minipage} \\ \hline
\end{tabularx}\medskip
\section*{Ejercicios}\vspace{-3mm}\color{LBcereza}\rule{\textwidth}{0.5pt}\color{LBtexto}
\needspace{5\baselineskip}
\par\medskip\noindent{\color{LBcian}\bfseries Consulta 1 -- Direcciones por distrito}\hfill{\color{LBcereza}\bfseries 2 puntos}\par\smallskip
Para cada distrito (\texttt{district}) registrado, muestra el nombre del distrito y el número de direcciones que hay en él. Ordena por cantidad descendente y, en caso de empate, por distrito en orden alfabético.\par\medskip
\needspace{5\baselineskip}
\par\medskip\noindent{\color{LBcian}\bfseries Consulta 2 -- Actores de películas largas}\hfill{\color{LBcereza}\bfseries 2 puntos}\par\smallskip
Para cada participación de un actor en una película, muestra el nombre y apellido del actor, el título de la película y la duración de esta en minutos, con el alias \texttt{duracion\_minutos}. Incluye únicamente las participaciones en películas de 150 minutos o más. Ordena por título ascendente y después por apellido del actor ascendente.\par\medskip
\needspace{5\baselineskip}
\par\medskip\noindent{\color{LBcian}\bfseries Consulta 3 -- Direcciones en India y China}\hfill{\color{LBcereza}\bfseries 2 puntos}\par\smallskip
Considera únicamente los países \texttt{India} y \texttt{China}. Para cada país, muestra su nombre, el número de ciudades diferentes que tienen al menos una dirección registrada y el número total de direcciones registradas. Ordena por total de direcciones descendente.\par\medskip
\newpage
\needspace{5\baselineskip}
\par\medskip\noindent{\color{LBcian}\bfseries Consulta 4 -- Películas PG-13 con reparto numeroso}\hfill{\color{LBcereza}\bfseries 2 puntos}\par\smallskip
Considera únicamente las películas con clasificación \texttt{PG-13}. Muestra el título de cada película y el número de actores de su reparto, pero solo para las que tengan 6 o más actores. Ordena por número de actores descendente y, en caso de empate, por título ascendente.\par\medskip
\needspace{5\baselineskip}
\par\medskip\noindent{\color{LBcian}\bfseries Consulta 5 -- Alquileres por día de la semana: clientes cuyo nombre empieza por S}\hfill{\color{LBcereza}\bfseries 2 puntos}\par\smallskip
Considera únicamente los alquileres de clientes cuyo nombre comienza por \texttt{S}. Para cada día de la semana de la fecha de alquiler, muestra el número de día de la semana (domingo = 1; sábado = 7), el número de alquileres y el número de esos alquileres que ya tienen fecha de devolución registrada. Ordena por número de día ascendente.\par\medskip
\medskip\noindent\begin{tabularx}{\textwidth}{|X|}\hline
\cellcolor{LBgris}\color{LBcian}\bfseries CRITERIOS DE CORRECCIÓN \\ \hline
\begin{minipage}{\dimexpr\linewidth-2\tabcolsep\relax}\vspace{2mm}\small Cada consulta vale 2 puntos. Se concederá puntuación parcial y no se penalizará dos veces un error arrastrado.\par\smallskip \begin{itemize}[leftmargin=4mm,nosep]\item \textbf{Datos y expresiones:} selección adecuada de columnas, cálculos y condiciones sencillas, así como los alias y el orden solicitado.
\item \textbf{Relación entre los datos:} utilización de las tablas y correspondencias necesarias para obtener los registros correctos, sin omisiones ni duplicaciones indebidas.
\item \textbf{Resumen de información:} obtención correcta de totales, cantidades, medias, grupos y condiciones sobre los resultados cuando sean necesarios.
\item \textbf{Corrección general:} resultado coherente y sintaxis razonablemente correcta. Se aceptarán formulaciones alternativas que resuelvan lo pedido.
\end{itemize}\vspace{2mm}\end{minipage} \\ \hline
\end{tabularx}\medskip

\end{document}
